Krátký glosář — rychlé definice a praktické tipy pro pedagogy. Uvedené záznamy kombinují stručné (general background knowledge) definice klíčových pojmů s konkrétními vzdělávacími příklady a odkazy na nástroje doporučené v této příručce. Cílem je poskytnout lehce použitelné „kapesní“ poznámky, které lze rychle adaptovat do výuky nebo pilotních aktivit.
\begin{itemize}
\item \textbf{Tokeny} (general background knowledge) --- základní jednotka textu, se kterou pracují velké jazykové modely; může to být část slova, celé slovo nebo interpunkce. (Použijte tuto intuici při tvorbě promptů a při odhadování délky požadavků na model.) Praktická rada: při zadávání úloh omezte délku vstupu, pokud chcete předvídatelnější výstupy.
\item \textbf{Embeddings} (general background knowledge) --- číselné (vektorové) reprezentace textu nebo dokumentů, užitečné pro vyhledávání podobných dokumentů, clustering nebo pro RAG‑workflow. Praktický nápad: předveďte studentům jednoduché hledání „nejpodobnějších“ odstavců z učebních materiálů.
\item \textbf{Halucinace (hallucination)} (general background knowledge) --- situace, kdy model generuje fakticky nesprávné nebo nepodložené informace; vždy ověřujte kritická fakta z nezávislých zdrojů a učte studenty metody verifikace. V hodině lze ukázat příklady správných a chybných odpovědí a diskutovat, jak chybu odhalit.
\item \textbf{Bias / předpojatost} (general background knowledge) --- systematické zkreslení výstupů modelu vyplývající z tréninkových dat nebo návrhu modelu; při výuce upozorněte studenty na možné zdroje zkreslení a zvažte kontrolní otázky. Aktivita: nechte studenty navrhnout prototesty, které odhalí možné zkreslení v konkrétních datech.
\item \textbf{RAG — Retrieval‑Augmented Generation} (general background knowledge) --- vzor, kdy se výstup modelu doplňuje o relevantní dokumenty získané z indexu/embeddingů před generováním; vhodné pro zajištění podložených odpovědí při práci s interními materiály. Ukažte rozdíl mezi odpovědí s a bez retrievalu na konkrétním odborném dotazu.
\item \textbf{Pokrok v matematice (Microsoft)} --- nástroj z rodiny „Vzdělávací akcelerátory“, který generuje sady příkladů, sbírá postupy studentů a nabízí předběžné opravy; užitečný příklad, jak AI v praxi zdůrazňuje sběr postupu nad pouhým výsledkem \cite{kb_ai_101_tipu_pdf_04e4a115_page_8_1}. Doporučení: použijte krátké cvičení na sběr postupu a následnou reflexi chyb.
\item \textbf{Minecraft Education} --- edukační prostředí s předpřipravenými světy (včetně aktivit „Hour of Code\"), vhodné pro názorné ukázky principů AI ve výuce a hands‑on demonstrace strojového učení \cite{kb_ai_101_tipu_pdf_04e4a115_page_15_2}. Praktické využití: 10min demonstrační aktivita ilustrující agentní chování nebo jednoduché senzory.
\item \textbf{Microsoft 365 Copilot — agent \"Výzkumník\"} --- doporučený pracovní režim Copilota pro hlubší rešerše a agregaci podkladů při přípravě komplexních výukových materiálů \cite{kb_ai_101_tipu_pdf_04e4a115_page_48_1}. Tip: použijte jej jako pomoc při sběru zdrojů, nikoli jako jediný ověřený zdroj.
\item \textbf{Interní repozitář VUT (ai.vut.cz)} --- living document s šablonami promptů, checklisty a video‑tutoriály; první místo, kam sáhnout po ověřených šablonách a krátkých návodech pro pilotní aktivity (kontaktujte metodické oddělení pro přístup a podporu). Obsah repozitáře pravidelně aktualizujte a sdílejte vlastní osvědčené postupy.
\end{itemize}
Rychlá praktická poznámka pro učitele: ke každému klíčovému pojmu mějte připravenou jednu krátkou ukázku/aktivitu (např. 10min demonstrace v Minecraft Education pro koncept AI; nebo generování dvou‑tří příkladů v Pokrok v matematice pro ilustraci významu sběru postupu) — tyto konkrétní nástroje jsou uvedeny výše a v repozitáři ai.vut.cz \cite{kb_ai_101_tipu_pdf_04e4a115_page_15_2,kb_ai_101_tipu_pdf_04e4a115_page_8_1,kb_ai_101_tipu_pdf_04e4a115_page_48_1}. Doporučujeme také krátké reflexní úlohy po každé aktivitě, které pomohou studentům kriticky zhodnotit výsledky.
Poznámka: stručné definice označené jako (general background knowledge) jsou záměrně minimální; pro výuku doporučujeme rozšířit je o krátké praktické ukázky a ověřovací úkoly z repozitáře. Přizpůsobte aktivity věkové skupině a cílovým kompetencím třídy a pravidelně iterujte podle zpětné vazby.